\chapter{Graph Neural Network Framework for NOMA User Pairing}
\label{chap:gnn}

The previous chapter established that traditional optimization methods face a fundamental trade-off between solution quality and computational complexity. Heuristic methods achieve $O(N \log N)$ complexity but sacrifice 20-30\% of potential throughput, while Blossom matching achieves global optimality but requires $O(N^3)$ complexity prohibitive for real-time operation. This chapter presents NOMA-GNN, our Graph Neural Network framework that bridges this gap by learning optimal pairing strategies offline, then executing real-time inference with near-linear complexity while maintaining 96.5\% of optimal performance.

\section{Problem Reformulation as Graph Learning}

We reformulate NOMA user pairing from discrete optimization to graph representation learning, explicitly modeling relationships between users and learning to predict optimal pairings from graph structure.

\subsection{Graph Construction}

Given $N$ users, we construct an undirected graph $\mathcal{G} = (\mathcal{V}, \mathcal{E})$ where:

\textbf{Nodes:} $\mathcal{V} = \{1, 2, \ldots, N\}$, one node per user

\textbf{Node Features:} Each node $i$ has a 5-dimensional feature vector capturing channel and spatial information:
\begin{equation}
\mathbf{x}_i = \begin{bmatrix} 
d_i \\ L_i \\ S_i \\ |h_i|^2 \\ h_i(\text{dB})
\end{bmatrix} \in \mathbb{R}^5
\end{equation}

where:
\begin{itemize}
    \item $d_i$ (m): Distance from base station to user $i$
    \item $L_i$ (dB): Path loss computed using 3GPP UMi model
    \item $S_i$ (dB): Log-normal shadowing component
    \item $|h_i|^2$: Squared magnitude of Rayleigh fading coefficient
    \item $h_i(\text{dB})$: Total channel gain in dB
\end{itemize}

These features capture complementary aspects of the wireless channel: distance-dependent path loss, shadowing due to obstacles, and small-scale multipath fading.

\textbf{Edges:} An edge $(i,j) \in \mathcal{E}$ exists if users $i$ and $j$ satisfy physical feasibility constraints:
\begin{equation}
(i,j) \in \mathcal{E} \iff 
\begin{cases}
|h_i - h_j|_{\text{dB}} \geq \Delta_{SIC} = 5 \text{ dB} & \text{(SIC constraint)} \\
|\theta_i - \theta_j| \leq \theta_{\max} = 90^\circ & \text{(Angular constraint)}
\end{cases}
\end{equation}

The SIC constraint ensures sufficient channel gain difference for successful interference cancellation. The angular constraint enhances spatial diversity by restricting pairings to users in similar angular sectors.

\textbf{Edge Features:} Each edge $(i,j)$ has features encoding pairability quality:
\begin{equation}
\mathbf{e}_{ij} = \begin{bmatrix} 
|h_i - h_j|_{\text{dB}} \\ |\theta_i - \theta_j| \\ d_{ij} \\ R_{ij}^{\text{est}}
\end{bmatrix} \in \mathbb{R}^4
\end{equation}

where $R_{ij}^{\text{est}}$ is an estimated sum-rate from a simple power allocation heuristic.

\subsection{Graph Properties}

Table~\ref{tab:graph_stats} shows typical graph statistics for different network sizes:

\begin{table}[ht]
\centering
\caption{Graph Statistics for Different Network Sizes}
\label{tab:graph_stats}
\begin{tabular}{|c|c|c|c|c|}
\hline
\textbf{Users $N$} & \textbf{Nodes $|\mathcal{V}|$} & \textbf{Edges $|\mathcal{E}|$} & \textbf{Avg Degree} & \textbf{Density} \\
\hline
20 & 20 & 89 & 8.9 & 0.47 \\
50 & 50 & 612 & 24.5 & 0.50 \\
100 & 100 & 2,458 & 49.2 & 0.50 \\
200 & 200 & 9,912 & 99.1 & 0.50 \\
\hline
\end{tabular}
\end{table}

The graph construction creates sparse structures where approximately 50\% of potential edges are feasible, significantly reducing the search space compared to fully connected graphs.

\section{GNN Architecture}

Our architecture follows an encoder-decoder paradigm: the encoder learns node representations from graph structure, while three specialized decoders predict pairing decisions, sum-rates, and power allocations.

\subsection{GraphSAGE Encoder}

We employ GraphSAGE (SAmple and AggreGatE) with 3 layers to capture multi-hop neighborhood information. Each layer performs:

\begin{equation}
\mathbf{h}_i^{(k)} = \text{ReLU}\left( \mathbf{W}^{(k)} \cdot \text{MEAN}\left(\{\mathbf{h}_j^{(k-1)} : j \in \mathcal{N}(i)\}\right) \right)
\end{equation}

where:
\begin{itemize}
    \item $\mathbf{h}_i^{(k)}$: Embedding of node $i$ at layer $k$
    \item $\mathbf{h}_i^{(0)} = \mathbf{x}_i$: Initial features
    \item $\mathcal{N}(i)$: Neighbors of node $i$
    \item $\text{MEAN}(\cdot)$: Mean aggregation
    \item $\mathbf{W}^{(k)}$: Learnable weight matrix
    \item $\text{ReLU}(\cdot)$: Rectified Linear Unit activation
\end{itemize}

\textbf{Layer Configuration:}
\begin{itemize}
    \item \textbf{Layer 1:} $\mathbb{R}^5 \rightarrow \mathbb{R}^{128}$ (input features to hidden)
    \item \textbf{Layer 2:} $\mathbb{R}^{128} \rightarrow \mathbb{R}^{128}$ (hidden to hidden)
    \item \textbf{Layer 3:} $\mathbb{R}^{128} \rightarrow \mathbb{R}^{128}$ (hidden to final)
    \item \textbf{Dropout:} 0.2 applied after each layer
\end{itemize}

After 3 GraphSAGE layers, node $i$ has representation $\mathbf{h}_i \in \mathbb{R}^{128}$ encoding information from its 3-hop neighborhood. This multi-hop aggregation is essential for capturing global pairing constraints—a user's pairing decision depends not only on its direct neighbors but also on what alternative pairings are available across the network.

\subsection{Multi-Task Decoder Architecture}

We employ three parallel decoders for joint prediction. For each edge $(i,j)$, we construct an edge representation by concatenating node embeddings and edge features:

\begin{equation}
\mathbf{z}_{ij} = [\mathbf{h}_i \,\|\, \mathbf{h}_j \,\|\, \mathbf{e}_{ij}] \in \mathbb{R}^{260}
\end{equation}

Each decoder is a 2-layer MLP that processes this edge representation:

\textbf{1. Pairing Decoder} (predicts probability edge belongs in matching):
\begin{align}
\mathbf{v}_{ij}^{(1)} &= \text{ReLU}(\mathbf{W}_{\text{pair}}^{(1)} \mathbf{z}_{ij} + \mathbf{b}_{\text{pair}}^{(1)}) \in \mathbb{R}^{128} \\
p_{ij} &= \sigma(\mathbf{W}_{\text{pair}}^{(2)} \mathbf{v}_{ij}^{(1)} + b_{\text{pair}}^{(2)}) \in [0,1]
\end{align}

where $\sigma(\cdot)$ is the sigmoid function and $p_{ij}$ represents the probability that users $i$ and $j$ should be paired.

\textbf{2. Sum-Rate Decoder} (predicts achievable sum-rate):
\begin{align}
\mathbf{v}_{ij}^{(2)} &= \text{ReLU}(\mathbf{W}_{\text{rate}}^{(1)} \mathbf{z}_{ij} + \mathbf{b}_{\text{rate}}^{(1)}) \in \mathbb{R}^{128} \\
\hat{R}_{ij} &= \mathbf{W}_{\text{rate}}^{(2)} \mathbf{v}_{ij}^{(2)} + b_{\text{rate}}^{(2)} \in \mathbb{R}^+
\end{align}

\textbf{3. Power Allocation Decoder} (predicts power allocation coefficient):
\begin{align}
\mathbf{v}_{ij}^{(3)} &= \text{ReLU}(\mathbf{W}_{\text{power}}^{(1)} \mathbf{z}_{ij} + \mathbf{b}_{\text{power}}^{(1)}) \in \mathbb{R}^{128} \\
\alpha_{ij} &= \sigma(\mathbf{W}_{\text{power}}^{(2)} \mathbf{v}_{ij}^{(3)} + b_{\text{power}}^{(2)}) \in [0,1]
\end{align}

where $\alpha_{ij}$ represents the fraction of total power allocated to the weaker user.

\subsection{Architecture Summary}

Algorithm~\ref{alg:gnn_forward} summarizes the complete forward pass:

\begin{algorithm}[H]
\caption{NOMA-GNN Forward Pass}
\label{alg:gnn_forward}
\begin{algorithmic}[1]
\Require Graph $\mathcal{G} = (\mathcal{V}, \mathcal{E})$, node features $\{\mathbf{x}_i\}$, edge features $\{\mathbf{e}_{ij}\}$
\Ensure Pairing probabilities $\{p_{ij}\}$, sum-rates $\{\hat{R}_{ij}\}$, power allocations $\{\alpha_{ij}\}$

\State \textbf{// Encoder: GraphSAGE Message Passing}
\State Initialize: $\mathbf{h}_i^{(0)} \leftarrow \mathbf{x}_i$ for all $i \in \mathcal{V}$
\For{$k = 1$ to $3$}
    \For{each node $i \in \mathcal{V}$}
        \State $\mathbf{a}_i \leftarrow \text{MEAN}\{\mathbf{h}_j^{(k-1)} : j \in \mathcal{N}(i)\}$
        \State $\mathbf{h}_i^{(k)} \leftarrow \text{ReLU}(\mathbf{W}^{(k)} \mathbf{a}_i)$
        \State $\mathbf{h}_i^{(k)} \leftarrow \text{Dropout}(\mathbf{h}_i^{(k)}, p=0.2)$
    \EndFor
\EndFor
\State Set final embeddings: $\mathbf{h}_i \leftarrow \mathbf{h}_i^{(3)}$ for all $i \in \mathcal{V}$

\State \textbf{// Decoder: Edge-Level Predictions}
\For{each edge $(i,j) \in \mathcal{E}$}
    \State $\mathbf{z}_{ij} \leftarrow [\mathbf{h}_i \,\|\, \mathbf{h}_j \,\|\, \mathbf{e}_{ij}]$
    \State $p_{ij} \leftarrow \text{PairingDecoder}(\mathbf{z}_{ij})$ \quad // 2-layer MLP with sigmoid
    \State $\hat{R}_{ij} \leftarrow \text{RateDecoder}(\mathbf{z}_{ij})$ \quad // 2-layer MLP
    \State $\alpha_{ij} \leftarrow \text{PowerDecoder}(\mathbf{z}_{ij})$ \quad // 2-layer MLP with sigmoid
\EndFor

\State \Return $\{p_{ij}\}, \{\hat{R}_{ij}\}, \{\alpha_{ij}\}$
\end{algorithmic}
\end{algorithm}

Table~\ref{tab:model_params} summarizes the model parameters:

\begin{table}[ht]
\centering
\caption{NOMA-GNN Model Parameters}
\label{tab:model_params}
\begin{tabular}{|l|l|r|}
\hline
\textbf{Component} & \textbf{Layer} & \textbf{Parameters} \\
\hline
\multirow{3}{*}{GraphSAGE Encoder} & Layer 1: $5 \rightarrow 128$ & 640 \\
& Layer 2: $128 \rightarrow 128$ & 16,384 \\
& Layer 3: $128 \rightarrow 128$ & 16,384 \\
\hline
\multirow{2}{*}{Pairing Decoder} & Hidden: $260 \rightarrow 128$ & 33,408 \\
& Output: $128 \rightarrow 1$ & 129 \\
\hline
\multirow{2}{*}{Sum-Rate Decoder} & Hidden: $260 \rightarrow 128$ & 33,408 \\
& Output: $128 \rightarrow 1$ & 129 \\
\hline
\multirow{2}{*}{Power Decoder} & Hidden: $260 \rightarrow 128$ & 33,408 \\
& Output: $128 \rightarrow 1$ & 129 \\
\hline
\textbf{Total} & & \textbf{166,275} \\
\hline
\end{tabular}
\end{table}

\section{Training Methodology}

\subsection{Training Data Generation}

We generate training data using the Blossom algorithm as an oracle. For each training sample:

\begin{enumerate}
    \item Generate random user deployment (positions, channel gains)
    \item Construct feasibility graph $\mathcal{G}$ with SIC/angular constraints
    \item Run Blossom algorithm to obtain optimal matching $\mathcal{M}^*$
    \item For each pair $(i,j) \in \mathcal{M}^*$, compute optimal power allocation and achievable sum-rate
    \item Create training labels: $y_{ij} = 1$ if $(i,j) \in \mathcal{M}^*$, else $y_{ij} = 0$
\end{enumerate}

\subsection{Multi-Task Loss Function}

We train with a composite loss combining three objectives:

\begin{equation}
\mathcal{L} = \mathcal{L}_{\text{pair}} + \lambda_R \mathcal{L}_{\text{rate}} + \lambda_\alpha \mathcal{L}_{\text{power}}
\end{equation}

where:

\textbf{1. Pairing Loss} (Binary Cross-Entropy):
\begin{equation}
\mathcal{L}_{\text{pair}} = -\frac{1}{|\mathcal{E}|} \sum_{(i,j) \in \mathcal{E}} \left[ y_{ij} \log p_{ij} + (1-y_{ij}) \log(1-p_{ij}) \right]
\end{equation}

\textbf{2. Sum-Rate Loss} (MSE on paired edges only):
\begin{equation}
\mathcal{L}_{\text{rate}} = \frac{1}{|\mathcal{M}^*|} \sum_{(i,j) \in \mathcal{M}^*} (\hat{R}_{ij} - R_{ij}^*)^2
\end{equation}

\textbf{3. Power Allocation Loss} (MSE on paired edges only):
\begin{equation}
\mathcal{L}_{\text{power}} = \frac{1}{|\mathcal{M}^*|} \sum_{(i,j) \in \mathcal{M}^*} (\alpha_{ij} - \alpha_{ij}^*)^2
\end{equation}

\textbf{Loss Weights:} $\lambda_R = 0.5$, $\lambda_\alpha = 0.5$

\subsection{Training Configuration}

Table~\ref{tab:training_config} shows the training hyperparameters:

\begin{table}[ht]
\centering
\caption{Training Configuration}
\label{tab:training_config}
\begin{tabular}{|l|l|}
\hline
\textbf{Hyperparameter} & \textbf{Value} \\
\hline
Optimizer & Adam \\
Learning Rate & $1 \times 10^{-3}$ \\
Weight Decay & $5 \times 10^{-4}$ \\
Batch Size & 2 \\
Number of Epochs & 60 \\
Dropout Rate & 0.2 \\
Early Stopping Patience & 10 epochs \\
\hline
\end{tabular}
\end{table}

\subsection{Training Algorithm}

Algorithm~\ref{alg:training} shows the training procedure:

\begin{algorithm}[H]
\caption{NOMA-GNN Training}
\label{alg:training}
\begin{algorithmic}[1]
\Require Training graphs $\{\mathcal{G}_t\}_{t=1}^T$ with optimal matchings $\{\mathcal{M}_t^*\}$
\Ensure Trained model parameters $\theta$

\State Initialize model parameters $\theta$ randomly
\State Initialize Adam optimizer with learning rate $\eta = 10^{-3}$
\State $\text{best\_loss} \leftarrow \infty$, $\text{patience} \leftarrow 0$

\For{epoch $e = 1$ to $60$}
    \State Shuffle training graphs
    \For{each batch of $B=2$ graphs}
        \State $\mathcal{L}_{\text{batch}} \leftarrow 0$
        \For{each graph $\mathcal{G}$ in batch}
            \State Forward pass: $\{p_{ij}\}, \{\hat{R}_{ij}\}, \{\alpha_{ij}\} \leftarrow \text{Forward}(\mathcal{G}; \theta)$
            \State Compute pairing loss: $\mathcal{L}_{\text{pair}}$
            \State Compute rate loss: $\mathcal{L}_{\text{rate}}$ (on paired edges only)
            \State Compute power loss: $\mathcal{L}_{\text{power}}$ (on paired edges only)
            \State $\mathcal{L} \leftarrow \mathcal{L}_{\text{pair}} + 0.5 \mathcal{L}_{\text{rate}} + 0.5 \mathcal{L}_{\text{power}}$
            \State $\mathcal{L}_{\text{batch}} \leftarrow \mathcal{L}_{\text{batch}} + \mathcal{L}$
        \EndFor
        \State Backward pass: compute $\nabla_\theta \mathcal{L}_{\text{batch}}$
        \State Update parameters: $\theta \leftarrow \theta - \eta \nabla_\theta \mathcal{L}_{\text{batch}}$
    \EndFor
    
    \State Evaluate on validation set: $\mathcal{L}_{\text{val}}$
    \If{$\mathcal{L}_{\text{val}} < \text{best\_loss}$}
        \State $\text{best\_loss} \leftarrow \mathcal{L}_{\text{val}}$
        \State Save model checkpoint
        \State $\text{patience} \leftarrow 0$
    \Else
        \State $\text{patience} \leftarrow \text{patience} + 1$
        \If{$\text{patience} \geq 10$}
            \State \textbf{break} \quad // Early stopping
        \EndIf
    \EndIf
\EndFor

\State \Return Best model parameters $\theta$
\end{algorithmic}
\end{algorithm}

\section{Inference and Matching Extraction}

At inference time, we need to convert edge probabilities $\{p_{ij}\}$ into a valid matching. We employ bipartite graph matching using NetworkX.

\subsection{Inference Algorithm}

Algorithm~\ref{alg:inference} shows the complete inference procedure:

\begin{algorithm}[H]
\caption{NOMA-GNN Inference}
\label{alg:inference}
\begin{algorithmic}[1]
\Require Graph $\mathcal{G} = (\mathcal{V}, \mathcal{E})$, trained model parameters $\theta$
\Ensure Matching $\mathcal{M}$, power allocations $\{\alpha_{ij} : (i,j) \in \mathcal{M}\}$

\State \textbf{// Step 1: GNN Forward Pass}
\State $\{p_{ij}\}, \{\hat{R}_{ij}\}, \{\alpha_{ij}\} \leftarrow \text{Forward}(\mathcal{G}; \theta)$

\State \textbf{// Step 2: Create Weighted Bipartite Graph}
\State Partition users: $\mathcal{V}_{\text{weak}} \leftarrow \{i : h_i < \text{median}\}$, $\mathcal{V}_{\text{strong}} \leftarrow \mathcal{V} \setminus \mathcal{V}_{\text{weak}}$
\State Create bipartite graph $G_B = (\mathcal{V}_{\text{weak}} \cup \mathcal{V}_{\text{strong}}, E_B)$
\State Set edge weights: $w_{ij} = -p_{ij}$ for all $(i,j)$ where $i \in \mathcal{V}_{\text{weak}}, j \in \mathcal{V}_{\text{strong}}$

\State \textbf{// Step 3: Maximum Weight Matching (NetworkX)}
\State $\mathcal{M} \leftarrow \text{nx.max\_weight\_matching}(G_B)$

\State \textbf{// Step 4: Extract Power Allocations}
\State $P \leftarrow \{\}$
\For{each $(i,j) \in \mathcal{M}$}
    \State $P[(i,j)] \leftarrow \alpha_{ij}$
\EndFor

\State \Return $\mathcal{M}, P$
\end{algorithmic}
\end{algorithm}

\subsection{Complexity Analysis}

Table~\ref{tab:complexity} compares computational complexity:

\begin{table}[ht]
\centering
\caption{Computational Complexity Comparison}
\label{tab:complexity}
\begin{tabular}{|l|c|c|}
\hline
\textbf{Method} & \textbf{Training} & \textbf{Inference} \\
\hline
Exhaustive Search & -- & $O(N!)$ \\
Blossom (Optimal) & -- & $O(N^3)$ \\
Static Heuristic & -- & $O(N \log N)$ \\
\hline
\textbf{NOMA-GNN} & $O(T \cdot |\mathcal{E}|)$ & $O(|\mathcal{E}| + N^3)$ \\
\quad (GNN forward) & & $O(|\mathcal{E}|) \approx O(N^2)$ \\
\quad (Matching extraction) & & $O(N^3)$ \\
\hline
\end{tabular}
\end{table}

\textbf{Note:} While the final matching step uses $O(N^3)$ bipartite matching, the GNN forward pass dominates practical runtime for large $N$ since $|\mathcal{E}| \ll N^2$ due to constraint-based edge pruning.

\subsection{Implementation Details}

The implementation uses the following software stack:

\begin{itemize}
    \item \textbf{Python 3.9:} Core programming language
    \item \textbf{PyTorch 1.13:} Neural network framework
    \item \textbf{PyTorch Geometric 2.2:} Graph neural network library
    \item \textbf{NetworkX 2.8:} Graph algorithms (bipartite matching)
    \item \textbf{NumPy 1.23:} Numerical computations
\end{itemize}

\section{Chapter Summary}

This chapter presented NOMA-GNN, a Graph Neural Network framework for learning optimal user pairing strategies:

\begin{itemize}
    \item \textbf{Graph Formulation:} Reformulated user pairing as graph learning with nodes representing users, edges representing feasible pairs, and labels indicating optimal matching decisions
    
    \item \textbf{GNN Architecture:} Employed 3-layer GraphSAGE encoder with 128 hidden dimensions to learn user embeddings, followed by three 2-layer MLP decoders for pairing, sum-rate, and power allocation predictions
    
    \item \textbf{Training:} Used Blossom-generated optimal matchings as supervision with multi-task loss combining binary cross-entropy (pairing), MSE (sum-rate), and MSE (power allocation)
    
    \item \textbf{Inference:} Converted GNN edge probabilities to valid matching using NetworkX bipartite matching algorithm
    
    \item \textbf{Model Size:} 166,275 trainable parameters (~650 KB checkpoint)
    
    \item \textbf{Complexity:} $O(|\mathcal{E}|)$ GNN forward pass + $O(N^3)$ matching extraction, competitive with traditional methods while maintaining near-optimal performance
\end{itemize}

The next chapter presents comprehensive experimental results demonstrating that NOMA-GNN achieves 96.5\% of optimal sum-rate while reducing inference time by 2-3 orders of magnitude compared to Blossom matching.
